\section{A Paired Transfer to Continuous-Time Welfare}
\label{sec:r15pairedtransfer}
\label{app:r15pairedtransfer}

The direct comparison of two executed policies permits a sharper transfer
than the sum of their separate discretization bounds. This section establishes
that transfer for the original capital diffusion and the original held-action
implementation. It changes neither the simulated observations nor the
confidence event. The construction uses only the economic primitives, the
action tube, and the assessment grid.

\subsection{The compared policies and the exact finite statistic}

Condition on all training and selection information. Let $a_k,b_k$ be the two
implemented actions on $[t_k,t_{k+1})$, where $t_k=kh$ and $Nh=T$. They may
depend on the entire innovation history through $t_k$, including the rounded
internal states used by their implementations. Both satisfy
\begin{equation}
 \|a_k-m_0(t_k)\bm1\|_\infty\leq\varepsilon,
 \qquad
 \|b_k-m_0(t_k)\bm1\|_\infty\leq\varepsilon.
 \label{eq:r15pairedtube}
\end{equation}
The inward guard in the executed verifier enforces these inequalities.
As in the capital implementation, this tube is contained in the primitive
positive consumption box.
The two policies receive the same initial profile and Brownian innovations.
Their physical states satisfy
\[
 dY_t^v=\{c\bm1+f(Y_t^v)-v_k\}\,dt+\Sigma\,dW_t,
 \qquad f(y)=\kappa\tanh(By),\qquad v\in\{a,b\}.
\]
Here $\Sigma=(\sigma_I I_d,\sigma_C\bm1)$ and
$c=\text{productivity}-(\sigma_I^2+\sigma_C^2)/2$.
The actions in these equations are the actions actually returned by the
original implementation. In particular, they are not recomputed at a
different state or on a different grid.

For the proof, introduce an auxiliary Euler state with exact arithmetic and
unclipped Brownian increments,
\begin{equation}
 Z_{k+1}^v=Z_k^v+h\{c\bm1+f(Z_k^v)-v_k\}
                 +\Sigma(W_{t_{k+1}}-W_{t_k}),
 \qquad Z_0^v=Y_0^v.
 \label{eq:r15pairedauxiliary}
\end{equation}
This is an analytical coupling. It is not a second executed policy. Its
actions remain those generated by the original clipped, rounded controller.
Consequently, no Lipschitz or differentiability assumption on the actor is
needed below.

Use $\langle x,y\rangle_d=d^{-1}x^\top y$,
$\|x\|_d^2=\langle x,x\rangle_d$, and
$\Pi=I_d-d^{-1}\bm1\bm1^\top$. Put
\[
 \psi(y)=\overline{f(y)},\qquad
 w(t)=\frac{1-e^{-\rho(T-t)}}{\rho}+e^{-\rho(T-t)},\qquad
 \omega(t)=e^{-\rho t}w(t),
\]
and define the exact scalar weights
\begin{equation}
 A_k=\int_{t_k}^{t_{k+1}}e^{-\rho t}\,dt,\qquad
 B_k=\int_{t_k}^{t_{k+1}}\omega(t)\,dt,\qquad
 J_k=\int_{t_k}^{t_{k+1}}\omega(t)(t-t_k)\,dt.
 \label{eq:r15pairedweights}
\end{equation}
Let $X_h^{a,b}$ be the following exact-real random variable:
\begin{align}
 X_h^{a,b}
 &=\sum_{k=0}^{N-1}\left[
 B_k\{\psi(Z_k^a)-\psi(Z_k^b)\}
 +A_k\{\overline{\log a_k}-\overline{\log b_k}\}
 -B_k(\bar a_k-\bar b_k)\right]\notag\\
 &\quad-\frac{\zeta}{2}\sum_{k=0}^{N-1}
       A_k(\bar a_k^2-\bar b_k^2)
 -\chi e^{-\rho T}
       \{\|\Pi Z_N^a\|_d^2-\|\Pi Z_N^b\|_d^2\}.
 \label{eq:r15pairedidealstatistic}
\end{align}
Integration by parts applied to
$e^{-\rho t}w(t)\overline{Y_t^v}$ shows that $J(a)-J(b)$ equals
the expectation of \eqref{eq:r15pairedidealstatistic}, with the production
sum replaced by
$\int_0^T\omega(t)\{\psi(Y_t^a)-\psi(Y_t^b)\}\,dt$
and the terminal states replaced by their physical counterparts.
The consumption terms in \eqref{eq:r15pairedidealstatistic} are already
integrated exactly because the implemented actions are held on each cell.
The common initial and exogenous-drift terms cancel.

\subsection{Computable global coefficients}

In this section $\|\cdot\|_2$ for a matrix is its Euclidean operator norm.
Choose a verified bound $b_*\geq\|B\|_2$ and set
\[
 \beta=\kappa b_*,\qquad
 L_2=\frac{4}{3\sqrt3},\qquad
 q_i=\sigma_I^2\|B_i\|_2^2+\sigma_C^2(B_i\bm1)^2,
 \qquad r_B=\max_i\|B_i\|_2.
\]
$B_i$ denotes row $i$ of $B$; the $q_i$ in this section are covariance
coefficients, not costates. Let
\[
 C_0=\max_{0\leq k<N}|c-m_0(t_k)|,\qquad
 M_\varepsilon=C_0+\kappa+\varepsilon,\qquad
 R_i=C_0|B_i\bm1|+(\kappa+\varepsilon)\sum_j|B_{ij}|.
\]
Any upper bounds for these quantities may be substituted. In the declared
economy, $m_0$ is increasing and $c<m_0(0)$, so the calculation bounds
$C_0$ by $m_0(T)-c$ using interval endpoints.

The scalar production map admits the Lipschitz bound
\begin{equation}
 |\psi(x)-\psi(y)|\leq\beta_\psi\|x-y\|_d,
 \quad
 \beta_\psi=
 \min\left\{\beta,
 \kappa\left[\frac1d\sum_{i,j}
                  \max\{(BB^\top)_{ij},0\}\right]^{1/2}\right\}.
 \label{eq:r15pairedscalargradient}
\end{equation}
Indeed, $dD\psi=\kappa B^\top u$ with $0\leq u_i\leq1$.
Bounding $u^\top BB^\top u$ either by $b_*^2d$ or by the sum of
the positive entries proves \eqref{eq:r15pairedscalargradient}.
The implementation replaces every displayed quantity by an outward upper
bound before taking the minimum.

Define
\begin{align}
 H_m&=\kappa L_2 b_*^2,&
 H_x&=\kappa L_2 b_*r_B\sqrt d,\notag\\
 G&=\kappa\left(\frac1d\sum_iq_i\right)^{1/2},&
 H_\Sigma&=\kappa L_2
       \|\operatorname{diag}(\sqrt q)B\|_2,\notag\\
 K&=\kappa\{C_0\|B\bm1\|_d+b_*(\kappa+\varepsilon)\}
                      +\frac{\kappa L_2}{2}\|q\|_d,\notag\\
 L_v&=\kappa L_2\|\operatorname{diag}(R)B\|_2
                         +\beta^2
                         +\kappa\|\operatorname{diag}(q)B\|_2,\notag\\
 L_m&=H_mM_\varepsilon+\beta_\psi\beta
       +\kappa\min\{b_*\|q\|_d,
                         \|\operatorname{diag}(q)B\|_2\}.
 \label{eq:r15pairedcoefficients}
\end{align}
All of these are functions of disclosed primitives and the action tube.
They do not depend on fitted weights, simulated outcomes, or an observed
maximum of a derivative.

To record the deterministic error recursion, let $r=2\varepsilon$ and put
\begin{align}
 D(t)&=r\frac{e^{\beta t}-1}{\beta},&
 I_1(t)&=\int_0^tD(s)\,ds,&
 I_2(t)&=\int_0^tD(s)^2\,ds,\notag\\
 \eta(t)&=he^{\beta t}
            \left\{\frac{Kt}{2}+G\sqrt{\frac t3}\right\}.
 \label{eq:r15pairedauxiliarybounds}
\end{align}
The continuous extension $D(t)=rt$ applies if $\beta=0$.
Set $\xi_0=0$ and, successively for $k=1,\ldots,N$, calculate
\begin{align}
 F_k&=\frac h2\{L_vI_1(t_k)+\beta r t_k\}
            +hH_\Sigma\sqrt{I_2(t_k)/3}
            +hH_x\sum_{j<k}D(t_j)\eta(t_j),\notag\\
 \xi_k&=F_k+h\beta\sum_{j<k}\xi_j.
 \label{eq:r15pairedrecursion}
\end{align}
This positive recursion retains the forcing accumulated through the current
date. It does not propagate a martingale through a future random Jacobian.

The initial distribution assigns probabilities $p_\ell$ to deterministic
profiles $y_\ell$. Define
\begin{equation}
 S_\ell=\|\Pi y_\ell\|_d+(\kappa+\varepsilon)T
              +\sigma_I\sqrt{(1-1/d)T},
 \qquad \bar S=\sum_\ell p_\ell S_\ell.
 \label{eq:r15pairedspread}
\end{equation}
The mean in this expression is an average of conditional bounds. It is not
asserted to be the $L^2$ norm of the state under the mixture of profiles.

\begin{proposition}[Paired payoff transfer]
\label{prop:r15pairedtransfer}
For the two held policies in \eqref{eq:r15pairedtube}, define
\begin{align}
 \mathfrak B_h^{\rm ideal}
 &=\sum_{k=0}^{N-1}B_k
          \{\beta_\psi\xi_k+H_mD(t_k)\eta(t_k)\}\notag\\
 &\quad+\sum_{k=0}^{N-1}J_k
       \left\{\frac{L_m}{h}[I_1(t_{k+1})-I_1(t_k)]
                                +\beta_\psi r\right\}\notag\\
 &\quad+\chi e^{-\rho T}
       \{(2\bar S+2D(T))\xi_N+2D(T)\eta(T)\}.
 \label{eq:r15pairedidealbound}
\end{align}
Then
\begin{equation}
 \left|J(a)-J(b)-\mathbb E X_h^{a,b}\right|
                       \leq\mathfrak B_h^{\rm ideal}.
 \label{eq:r15pairedidealclaim}
\end{equation}
Suppose $\delta_{\rm num}$ bounds, in conditional $L^2$ at every node,
the distance between each auxiliary Euler state and its original rounded,
clipped internal state. Suppose $\epsilon_{\rm stat}$ bounds the complete
arithmetic error of each stored policy-minus-reference statistic, and
$\epsilon_{\rm sub}$ bounds their final subtraction. Then the original
stored direct difference $\widehat X_h^{a,b}$ satisfies
\begin{equation}
 \left|J(a)-J(b)-\mathbb E\widehat X_h^{a,b}\right|
                       \leq\mathfrak B_h^{\rm pair},
 \label{eq:r15pairedimplementedclaim}
\end{equation}
where
\begin{align}
 \mathfrak B_h^{\rm pair}
 &=\mathfrak B_h^{\rm ideal}
       +2\beta_\psi\delta_{\rm num}\sum_kB_k\notag\\
 &\quad+2\chi e^{-\rho T}\delta_{\rm num}
                           (2\bar S+\delta_{\rm num})
       +2\epsilon_{\rm stat}+\epsilon_{\rm sub}.
 \label{eq:r15pairedimplementedbound}
\end{align}
The conclusion is uniform over both policies satisfying the displayed tube.
\end{proposition}

\subsection{Proof of Proposition~\ref{prop:r15pairedtransfer}}

\paragraph*{Regularity and integrability.}
The drift is globally Lipschitz in the state, with constant $\beta$, and
bounded uniformly over admissible actions. Thus the physical equations have
unique strong solutions with finite moments on the finite horizon. The
actions are bounded and measurable at the left endpoint of their cell.
The first three derivatives of $f$ are bounded. All stochastic integrals
used below are consequently square integrable, so their expectations vanish
and It\^o's isometry applies without an unaccounted localization limit.
The quadratic terminal payoff and every product below are integrable.

\paragraph*{Derivative and generator bounds.}
The elementary identities for $\tanh$ give
$|\tanh'|\leq1$, $|\tanh''|\leq L_2$, and
$|\tanh'''|\leq2$. For a held action $v$, write
\[
 \mathcal L^v u(y)
 =Du(y)\{c\bm1+f(y)-v\}
            +\tfrac12\operatorname{tr}\{\Sigma\Sigma^\top D^2u(y)\}.
\]
For vector-valued $u$, this definition applies to each component.
Since $v=m_0(t_k)\bm1+e$ with $\|e\|_\infty\leq\varepsilon$,
\[
 \|\mathcal L^v f(y)\|_d\leq K,
 \qquad d^{-1/2}\|Df(y)\Sigma\|_{\rm HS}\leq G.
\]
To see the first inequality, bound
$\|Df(y)(c\bm1+f(y)-v)\|_d$ by
$\kappa\{C_0\|B\bm1\|_d+b_*(\kappa+\varepsilon)\}$.
The $i$th second-order term is
$\kappa q_i\tanh''(B_i y)/2$, whose normalized norm is at most
$\kappa L_2\|q\|_d/2$.

For two states $x,y$ and two actions $v,w$ in the same tube, splitting
the drift part into the change in $Df$ and the change in drift gives
\begin{align}
 \|\mathcal L^v f(x)-\mathcal L^w f(y)\|_d
          &\leq L_v\|x-y\|_d+\beta\|v-w\|_d,\notag\\
 |\mathcal L^v\psi(x)-\mathcal L^w\psi(y)|
          &\leq L_m\|x-y\|_d+\beta_\psi\|v-w\|_d.
 \label{eq:r15pairedgeneratorbounds}
\end{align}
In the first line, the change in $Df$ is bounded by
$\kappa L_2\|\operatorname{diag}(R)B\|_2\|x-y\|_d$,
because $|B_i(c\bm1+f(x)-v)|\leq R_i$.
The remaining first-order terms are bounded by
$\beta^2\|x-y\|_d+\beta\|v-w\|_d$.
The change in the second-order term is bounded by
$\kappa\|\operatorname{diag}(q)B\|_2\|x-y\|_d$.

For the second line, the scalar Hessian obeys
$\|D^2\psi\|_2\leq H_m/d$. Its contribution to the drift difference
is at most $H_mM_\varepsilon\|x-y\|_d$; the remaining drift terms
contribute $\beta_\psi\beta\|x-y\|_d+
\beta_\psi\|v-w\|_d$.
The change in its second-order term is bounded by either
$\kappa b_*\|q\|_d\|x-y\|_d$ or
$\kappa\|\operatorname{diag}(q)B\|_2\|x-y\|_d$.
This proves the coefficients in \eqref{eq:r15pairedcoefficients}.

The normalized Hilbert--Schmidt bound for the paired diffusion derivative is
\begin{align}
 \frac1d\|\{Df(x)-Df(y)\}\Sigma\|_{\rm HS}^2
 &=\frac{\kappa^2}{d}\sum_i q_i
        \{\tanh'(B_i x)-\tanh'(B_i y)\}^2\notag\\
 &\leq\frac{\kappa^2L_2^2}{d}
            \|\operatorname{diag}(\sqrt q)B(x-y)\|_2^2
 \leq H_\Sigma^2\|x-y\|_d^2.
 \label{eq:r15pairedhsmatrix}
\end{align}
This calculation accounts for the dimension normalization explicitly.
The scalar production Hessian likewise has no extra $\sqrt d$ factor.
For the vector four-point bound, an operator-norm estimate does require
\[
 \|Df(x)-Df(y)\|_2
 \leq\kappa L_2 b_*\max_i|B_i(x-y)|
 \leq H_x\|x-y\|_d.
\]

\paragraph*{The individual and paired state errors.}
Fix an initial profile. Let
$\Delta Y_t=Y_t^a-Y_t^b$, $\Delta Z_k=Z_k^a-Z_k^b$, and
$e_k^v=Y_{t_k}^v-Z_k^v$.
The common Brownian increments cancel from each state difference. Since
$\|a_k-b_k\|_d\leq r$, pathwise Gronwall gives
\begin{equation}
 \|\Delta Y_t\|_d\leq D(t),\qquad
 \|\Delta Z_k\|_d\leq
       r\frac{(1+h\beta)^k-1}{\beta}\leq D(t_k).
 \label{eq:r15paireddistance}
\end{equation}
For a single policy the accumulated local drift error through $t_k$ is
\[
 R_k^v=\sum_{j<k}\int_{t_j}^{t_{j+1}}
                 \{f(Y_s^v)-f(Y_{t_j}^v)\}\,ds.
\]
Apply It\^o's formula to $f(Y_s^v)$ within each cell and interchange
the integrals. If $\theta_h(s)=t_{j+1}-s$ for
$s\in[t_j,t_{j+1})$, Minkowski's inequality and It\^o's isometry yield
\[
 \|R_k^v\|_{L^2(d)}
 \leq K\int_0^{t_k}\theta_h(s)\,ds
       +G\left\{\int_0^{t_k}\theta_h(s)^2\,ds\right\}^{1/2}
 =h\{Kt_k/2+G\sqrt{t_k/3}\}.
\]
Here $\|U\|_{L^2(d)}=(\mathbb E\|U\|_d^2)^{1/2}$, with
expectation conditional on the fixed initial profile. The nodal error
recursion and discrete Gronwall therefore give
\begin{equation}
 \|e_k^v\|_{L^2(d)}\leq\eta(t_k).
 \label{eq:r15pairedsingleerror}
\end{equation}

Apply the same representation to $R_k^a-R_k^b$.
Equations \eqref{eq:r15pairedgeneratorbounds} and
\eqref{eq:r15pairedhsmatrix} imply
\begin{align*}
 \|R_k^a-R_k^b\|_{L^2(d)}
 &\leq\int_0^{t_k}\theta_h(s)\{L_vD(s)+\beta r\}\,ds\\
 &\quad+H_\Sigma
       \left\{\int_0^{t_k}\theta_h(s)^2D(s)^2\,ds\right\}^{1/2}\\
 &\leq\frac h2\{L_vI_1(t_k)+\beta r t_k\}
                         +hH_\Sigma\sqrt{I_2(t_k)/3}.
\end{align*}
For the last inequality, on each cell $\theta_h$ and
$\theta_h^2$ decrease while $D$ and $D^2$ increase. The integral
of a decreasing function times an increasing function is no greater than
the product of their cell averages times the cell length. Their kernel
averages are $h/2$ and $h^2/3$, respectively.
The martingale isometry is applied to the single accumulated stochastic
integral with predictable integrand. There is no isometry assertion for
a future random flow multiplying a past noise increment.

For the difference of the nodal drift errors, the four-point identity and
the derivative bound give
\begin{align*}
 &\|f(Y_{t_j}^a)-f(Y_{t_j}^b)
                         -f(Z_j^a)+f(Z_j^b)\|_d\\
 &\qquad\leq\beta\|e_j^a-e_j^b\|_d
                    +H_x\|\Delta Z_j\|_d\|e_j^b\|_d.
\end{align*}
One obtains this inequality by first changing $\Delta Y_{t_j}$ to
$\Delta Z_j$ at the base point $Y_{t_j}^b$, and then changing that
base point to $Z_j^b$ with $\Delta Z_j$ fixed.
Combining it with \eqref{eq:r15paireddistance} and
\eqref{eq:r15pairedsingleerror} proves inductively that
\begin{equation}
 \|e_k^a-e_k^b\|_{L^2(d)}\leq\xi_k.
 \label{eq:r15pairednodalerror}
\end{equation}
All bounds are uniform in the initial profile.

\paragraph*{Production at and between decision dates.}
The corresponding scalar four-point inequality uses $H_m$ in place of
$H_x$ and $\beta_\psi$ in place of $\beta$. Thus
\[
 \mathbb E|\psi(Y_{t_k}^a)-\psi(Y_{t_k}^b)
                       -\psi(Z_k^a)+\psi(Z_k^b)|
 \leq\beta_\psi\xi_k+H_mD(t_k)\eta(t_k).
\]
Multiplication by $B_k$ accounts for the first line of
\eqref{eq:r15pairedidealbound}.

For variation inside a cell, take expectations in It\^o's formula for
$\psi(Y_t^a)-\psi(Y_t^b)$. The martingale has mean zero by the
integrability already established. Put
$\Psi_k(s)=\int_s^{t_{k+1}}\omega(t)\,dt$. The expected weighted
within-cell error is bounded in absolute value by
\[
 \int_{t_k}^{t_{k+1}}\Psi_k(s)
                          \{L_mD(s)+\beta_\psi r\}\,ds
 \leq J_k\left\{\frac{L_m}{h}
       [I_1(t_{k+1})-I_1(t_k)]+\beta_\psi r\right\}.
\]
Here $\Psi_k$ decreases, $D$ increases, and
$\int_{t_k}^{t_{k+1}}\Psi_k(s)\,ds=J_k$.
This proves the second line of \eqref{eq:r15pairedidealbound}.

\paragraph*{Terminal dispersion and the initial population.}
For a fixed profile $y_\ell$, projecting the physical state equation removes
the scalar schedule, the exogenous scalar drift, and the common shock.
Minkowski's inequality yields
$\|\Pi Y_T^b\|_{L^2(d)}\leq S_\ell$.
The same bound holds for the auxiliary Euler state $\Pi Z_N^b$ and for
either policy separately. Use the exact terminal identity
\begin{align*}
 &\|\Pi Y_T^a\|_d^2-\|\Pi Y_T^b\|_d^2
                 -\|\Pi Z_N^a\|_d^2+\|\Pi Z_N^b\|_d^2\\
 &=2\langle\Pi Y_T^b,\Pi(e_N^a-e_N^b)\rangle_d
       +2\langle\Pi e_N^b,\Pi\Delta Z_N\rangle_d\\
 &\quad+\langle\Pi(\Delta Y_T+\Delta Z_N),
                              \Pi(e_N^a-e_N^b)\rangle_d.
\end{align*}
Conditional Cauchy--Schwarz and the pathwise bounds for both differences
give the upper bound
$(2S_\ell+2D(T))\xi_N+2D(T)\eta(T)$ on the absolute expectation.
Average this bound over the profiles. Because $D,\eta,\xi$ are uniform
in the initial state, the result contains $\bar S=\sum p_\ell S_\ell$.
In particular, the calculation uses the mean initial spread after applying
Cauchy--Schwarz separately for each profile; it does not replace a mixture
second moment by a mean of square roots. Multiplication by
$\chi e^{-\rho T}$ proves the remaining line of
\eqref{eq:r15pairedidealbound}, and hence
\eqref{eq:r15pairedidealclaim}.

\paragraph*{The stored statistic.}
Let $\widetilde Z_k^v$ denote the actual rounded, clipped internal state.
The inherited forward-arithmetic and normal-clipping bounds, with the
actions held fixed as above, provide a uniform value
$\delta_{\rm num}\geq
\|Z_k^v-\widetilde Z_k^v\|_{L^2(d)}$ conditional on each profile.
Scalar production then contributes at most
$2\beta_\psi\delta_{\rm num}\sum B_k$ for the two policies.
For each policy,
\[
 \mathbb E\left|\|\Pi Z_N^v\|_d^2
                         -\|\Pi\widetilde Z_N^v\|_d^2\right|
 \leq\delta_{\rm num}(2S_\ell+\delta_{\rm num}).
\]
Averaging over the profiles and adding the two bounds gives the terminal
arithmetic term in \eqref{eq:r15pairedimplementedbound}.

Each original stored gain is a policy-minus-reference statistic. Under the
verified common-noise and initial-state identities, their exact-real
reference terms coincide, and cancel in the direct difference. Charging
the complete stored-statistic error for each policy is sufficient even
though some numerical terms also cancel. The final floating-point
subtraction receives its original additional allowance. This proves
\eqref{eq:r15pairedimplementedclaim}.\hfill Q.E.D.

\subsection{Outward evaluation and reuse of the original confidence event}

The numerical account uses the original $N=2048$ grid, $T=1$, and
$\varepsilon=0.1$. Every coefficient is evaluated with the inherited
interval arithmetic. The coupling norm is accepted only after positive
interval LDL pivots certify a matrix upper bound. For an interval-valued
weighted matrix $C$, the computation uses
\[
 \|C\|_2\leq\|C_{\rm mid}\|_2+\|C-C_{\rm mid}\|_{\rm F},
\]
with a verified norm for the binary64 midpoint and an outward Frobenius
bound for its interval radius. Componentwise upper bounds for $q$ and
$R$ are legitimate in the weighted norms, since increasing nonnegative
diagonal weights can only increase the Euclidean norm of each image.

Positive power series avoid cancellation in $D$, $I_1$, and $I_2$ near
the first node. They are
\begin{align*}
 D(t)&=r\sum_{m\geq0}\frac{\beta^m t^{m+1}}{(m+1)!},\\
 I_1(t)&=r\sum_{m\geq0}\frac{\beta^m t^{m+2}}{(m+2)!},\\
 I_2(t)&=r^2\sum_{m\geq0}
           \frac{(2^{m+2}-2)\beta^m t^{m+3}}{(m+3)!}.
\end{align*}
The implementation retains $m=0,\ldots,28$.
Writing $z_*=\beta T<1$, the successive ratios in each remaining tail
are bounded by
$z_*/30$, $z_*/31$, and $3z_*/32$, respectively.
If $u$ is the last retained term and $r_*<1$ is its tail-ratio bound,
the omitted positive tail is at most $ur_*/(1-r_*)$.
All sums, products, roots, differences of integral enclosures, and the
positive recursion \eqref{eq:r15pairedrecursion} are rounded outward.
The scalar $B_k,J_k$ are their original interval enclosures.

For the implementation terms, the calculation takes the maximum, over the
nine initial profiles, of the inherited forward-roundoff plus
normal-clipping strong-error bound. It also retains twice the complete
original stored-statistic arithmetic allowance and the original
$10^{-12}$ allowance for the direct subtraction. The latter includes
no proposed cancellation of scalar quadrature or evaluation errors.
The calculation is conditional on the same binary64 and ideal-iid-normal
sampling contract as the original verifier.

This account applies when both compared outputs are held controllers.
If both outputs select the analytical schedule, their difference remains
exactly zero. If exactly one selects that schedule, the original complete
schedule-relative transfer is retained. A continuously varying analytical
schedule is not silently treated as a held controller.

Finally, let the original empirical-Bernstein event bound the mean of the
same clipped direct-difference observations, with its original clipping
threshold and tail allowance. Both the original deterministic transfer
and \eqref{eq:r15pairedimplementedclaim} hold uniformly for the two
held controllers. Their minimum may therefore be used on that same
event. This operation spends no additional probability budget, changes
no sample or stopping rule, and makes no claim of a smaller empirical
variance. The original endpoints remain available together with the
refined endpoints and the deterministic accounts that distinguish them.
